\documentclass[../Odyssey_Project.tex]{subfiles}

\begin{document}
\setcounter{section}{5}
\section{The Special Linear Group}

\subsubsection*{Definition}
\begin{itemize}
    \item The \emph{special linear group} is the subgroup $\mathrm{SL}(n,F)$ of $\mathrm{GL}(n,F)$ consisting of all matrices having determinant $1$.
    \item In other words, $\mathrm{SL}(n,F)$ is the kernel of the homomorphism $\det\colon \mathrm{GL}(n,F) \to F^{\times}$, and hence $\mathrm{SL}(n,F) \trianglelefteq \mathrm{GL}(n,F)$.
\end{itemize}

\begin{proposition}
\label{prop:6.1}
Let $n \in \N$, and let $q$ be a prime power. Then
\[
|\mathrm{SL}(n,q)| = \prod_{k=1}^{n-1}(q^{n+1}-q^k) = q^{\frac{n(n-1)}{2}}(q^n-1)\cdots(q^2-1).
\]
\end{proposition}
\begin{proof}
Let $F$ be the field with $q$ elements. The determinant map
\[
    \det\colon \GL(n,F)\to F^\times
\]
is surjective: if $\lambda\in F^\times$, then the diagonal matrix with first diagonal entry $\lambda$ and all other diagonal entries equal to $1$ has determinant $\lambda$. Its kernel is precisely $\SL(n,F)$. Hence, by the fundamental theorem on homomorphisms,
\[
    \GL(n,F)/\SL(n,F)\cong F^\times.
\]
In particular,
\[
    \lvert \GL(n,q) : \SL(n,q) \rvert=\lvert F^\times\rvert=q-1.
\]
Proposition~\ref{prop:4.2} gives
\[
    \lvert \GL(n,q)\rvert
    =q^{\frac{n(n-1)}{2}}(q^n-1)(q^{n-1}-1)\cdots(q-1).
\]
Dividing by the index $q-1$ cancels the final factor, so
\[
    \lvert \SL(n,q)\rvert
    =\frac{\lvert \GL(n,q)\rvert}{q-1}
    =q^{\frac{n(n-1)}{2}}(q^n-1)\cdots(q^2-1).
\]
Finally,
\[
    \prod_{k=1}^{n-1}(q^{n+1}-q^k)
    =\prod_{k=1}^{n-1}q^k(q^{n+1-k}-1)
    =q^{1+\cdots+(n-1)}(q^n-1)\cdots(q^2-1),
\]
which is the same expression.
\end{proof}

\subsubsection*{Generation by root subgroups}
\begin{itemize}
    \item We showed in Theorem~\ref{thm:4.12} that $\mathrm{GL}(n,F)$ is generated by the root subgroups and the diagonal subgroup. We now establish the analogous result for $\mathrm{SL}(n,F)$.
\end{itemize}

\begin{proposition}
\label{prop:6.2}
$\mathrm{SL}(n,F)$ is generated by the root subgroups $X_{ij}$.
\end{proposition}
\begin{proof}
If $n=1$, then $\SL(1,F)=1$, so there is nothing to prove. We assume now that $n\geq 2$. Let $H$ be the subgroup of $\GL(n,F)$ generated by all root subgroups $X_{ij}$. By Lemma~\ref{lem:4.6}(i), every transvection has determinant $1$, so $H\leq \SL(n,F)$. We prove the reverse containment.\\
Let $A\in \SL(n,F)$. We shall premultiply $A$ by products of transvections until it becomes the identity. Since the inverse of a transvection is again a transvection by Lemma~\ref{lem:4.6}(iii), this will imply that $A\in H$.\\
First we make $A$ upper triangular. By Lemma~\ref{lem:4.7}, there is a product $b$ of upper triangular transvections such that, for each $1\leq i\leq n$, the matrix $bA$ has exactly one row whose entries in the first $i-1$ columns are zero and whose entry in the $i$th column is non-zero. In other words, the rows of $bA$ already have the correct pivot pattern for an upper triangular matrix, but the pivot rows may be in the wrong order.\\
For distinct $i$ and $j$, put
\[
    s_{ij}=X_{ji}(1)X_{ij}(-1)X_{ji}(1).
\]
On the two rows indexed by $i$ and $j$, premultiplication by $s_{ij}$ has the same effect as the matrix
\[
    \begin{pmatrix}
        0 & -1\\
        1 & 0
    \end{pmatrix};
\]
that is, it sends row $i$ to the negative of row $j$ and sends row $j$ to row $i$. Thus $s_{ij}$ swaps the two rows up to a harmless sign. Using such signed row swaps, we can move the row with the first pivot into the first row, the row with the second pivot into the second row, and so on. Therefore there is a product $c$ of transvections such that
\[
    cbA
\]
is upper triangular.\\
Next we make the diagonal entries equal to $1$. Let $M$ be an upper triangular element of $\SL(n,F)$, with diagonal entries $\lambda_1,\ldots,\lambda_n$. We use the following two-by-two calculation. If $a,c\in F^\times$ and $b'\in F$, then
\[
    X_{21}(-1)X_{12}(1-c^{-1})X_{21}(c)
    =
    \begin{pmatrix}
        c & 1-c^{-1}\\
        0 & c^{-1}
    \end{pmatrix},
\]
and hence
\[
    X_{21}(-1)X_{12}(1-c^{-1})X_{21}(c)
    \begin{pmatrix}
        a & b'\\
        0 & c
    \end{pmatrix}
    =
    \begin{pmatrix}
        ac & cb'+c-1\\
        0  & 1
    \end{pmatrix},
\]
so the only point is that the two diagonal entries $a,c$ have become $ac,1$. Applying this calculation to the final two diagonal positions of $M$ changes the final two diagonal entries from $\lambda_{n-1},\lambda_n$ to $\lambda_{n-1}\lambda_n,1$, while keeping the matrix upper triangular. Applying the same calculation next to positions $n-2,n-1$, then to positions $n-3,n-2$, and continuing leftward, we obtain an upper triangular matrix whose diagonal entries are
\[
    \lambda_1\lambda_2\cdots\lambda_n,1,\ldots,1.
\]
Since $M\in \SL(n,F)$, the product of the diagonal entries of $M$ is its determinant, and hence $\lambda_1\cdots\lambda_n=1$. Thus a product of transvections premultiplies $M$ to an upper unitriangular matrix. Applying this to $M=cbA$, we find a product $d$ of transvections such that $dcbA$ is upper unitriangular.\\
Finally, by Lemma~\ref{lem:4.11}, an upper triangular matrix can be premultiplied by a product of transvections to become the diagonal matrix with the same main diagonal. Since $dcbA$ is upper unitriangular, this diagonal matrix is the identity. Hence there is a product $e$ of transvections such that
\[
    edcbA=I.
\]
It follows that
\[
    A=b^{-1}c^{-1}d^{-1}e^{-1}.
\]
Each factor on the right lies in $H$, so $A\in H$. Therefore $\SL(n,F)\leq H$, and hence $\SL(n,F)=H$.
\end{proof}

\begin{proposition}
\label{prop:6.3}
The subgroups $X_{ij}$ are conjugate in $\mathrm{SL}(n,F)$.
\end{proposition}
\begin{proof}
If $n=1$, there are no root subgroups, so there is nothing to prove. Let $X_{ij}$ and $X_{i'j'}$ be two root subgroups. Choose a permutation matrix $w$ such that $w(i)=i'$ and $w(j)=j'$. By Lemma~\ref{lem:4.6}(v), for every $\alpha\in F$ we have
\[
    wX_{ij}(\alpha)w^{-1}=X_{i'j'}(\alpha).
\]
If $\det w=1$, then $w\in \SL(n,F)$, and hence $wX_{ij}w^{-1}=X_{i'j'}$.\\
It remains to handle the case $\det w=-1$. Let $d$ be the diagonal matrix whose first diagonal entry is $-1$ and whose remaining diagonal entries are $1$. Then $\det d=-1$, so $\det(dw)=1$, and hence $dw\in \SL(n,F)$. For each $\alpha\in F$,
\[
    (dw)X_{ij}(\alpha)(dw)^{-1}
    =dX_{i'j'}(\alpha)d^{-1}.
\]
Conjugating by $d$ can only change the parameter $\alpha$ by a sign, depending on whether one of $i'$ or $j'$ is equal to $1$. Thus
\[
    dX_{i'j'}(\alpha)d^{-1}=X_{i'j'}(\pm\alpha).
\]
As $\alpha$ ranges over $F$, so does $\pm\alpha$. Therefore
\[
    (dw)X_{ij}(dw)^{-1}=X_{i'j'}.
\]
Thus any two root subgroups are conjugate in $\SL(n,F)$.
\end{proof}

\subsubsection*{Scalar matrices and the center}
\begin{itemize}
    \item Let $Z$ be the subgroup of $\mathrm{GL}(n,F)$ consisting of the multiples of the identity matrix by elements of $F^{\times}$. The elements of $Z$ are called the (non-zero) \emph{scalar} matrices. We clearly have $Z \cong F^{\times}$.
\end{itemize}

\begin{proposition}
\label{prop:6.4}
The center of $\mathrm{GL}(n,F)$ is $Z$, and the center of $\mathrm{SL}(n,F)$ is $Z \cap \mathrm{SL}(n,F)$.
\end{proposition}
\begin{proof}
If $n=1$, then $\GL(1,F)=F^\times=Z$ and $\SL(1,F)=1=Z\cap \SL(1,F)$, so the result is immediate. We assume now that $n>1$. Let $G$ denote either $\GL(n,F)$ or $\SL(n,F)$. Since scalar matrices commute with every matrix, we have
\[
    Z\cap G\leq Z(G).
\]
It remains to prove the reverse containment. It is enough to show that if an element of $G$ commutes with every element of $\SL(n,F)$, then it is scalar.\\
Let $M=(m_{rs})$ be such an element of $G$. Fix distinct indices $i$ and $j$. By Lemma~\ref{lem:4.6}(i), the transvection $X_{ij}(1)$ lies in $\SL(n,F)$, so
\[
    MX_{ij}(1)=X_{ij}(1)M.
\]
Writing $X_{ij}(1)=I+E_{ij}$, this equality is equivalent to
\[
    ME_{ij}=E_{ij}M.
\]
Compare entries. The $(i,i)$-entry of $ME_{ij}$ is $0$, while the $(i,i)$-entry of $E_{ij}M$ is $m_{ji}$, so $m_{ji}=0$. The $(i,j)$-entry of $ME_{ij}$ is $m_{ii}$, while the $(i,j)$-entry of $E_{ij}M$ is $m_{jj}$, so $m_{ii}=m_{jj}$.\\
Since $i$ and $j$ were arbitrary distinct indices, every off-diagonal entry of $M$ is zero, and all diagonal entries are equal. Hence $M=m_{11}I$ is a scalar matrix. Thus $M\in Z\cap G$. Therefore $Z(G)=Z\cap G$. Taking $G=\GL(n,F)$ gives $Z(\GL(n,F))=Z$, and taking $G=\SL(n,F)$ gives $Z(\SL(n,F))=Z\cap \SL(n,F)$.
\end{proof}

\subsubsection*{Projective groups}
\begin{itemize}
    \item The group $\mathrm{GL}(n,F)/Z$ is called the \emph{projective general linear group}, and the group $\mathrm{SL}(n,F)/Z \cap \mathrm{SL}(n,F)$ is called the \emph{projective special linear group}.
    \item We denote these groups by $\mathrm{PGL}(n,F)$ and $\mathrm{PSL}(n,F)$, respectively. (Observe that the first isomorphism theorem implies that $\mathrm{PSL}(n,F) \cong Z\,\mathrm{SL}(n,F)/Z \leq \mathrm{GL}(n,F)/Z = \mathrm{PGL}(n,F)$.)
    \item As $Z \cong F^{\times}$, we have $|\mathrm{PGL}(n,q)| = |\mathrm{SL}(n,q)|$ for any prime power $q$; however, to determine $|\mathrm{PSL}(n,q)|$ we must be able to count the number of $n$th roots of unity in the field of $q$ elements.
\end{itemize}

\begin{proposition}
\label{prop:6.5}
Let $q$ be a prime power. Then
\[
|\mathrm{PSL}(2,q)| = \begin{cases} q^3-q & \text{if } 2 \mid q, \\[2pt] \tfrac{1}{2}(q^3-q) & \text{if } 2 \nmid q. \end{cases}
\]
\end{proposition}
\begin{proof}
Let $F$ be the field with $q$ elements. By Proposition~\ref{prop:6.4},
\[
    Z(\SL(2,F))=Z\cap \SL(2,F),
\]
and hence
\[
    \lvert \mathrm{PSL}(2,q)\rvert
    =\frac{\lvert \SL(2,q)\rvert}{\lvert Z\cap \SL(2,q)\rvert}.
\]
By Proposition~\ref{prop:6.1}, we have
\[
    \lvert \SL(2,q)\rvert=q(q^2-1)=q^3-q.
\]
It remains to count the scalar matrices in $\SL(2,q)$. Such a matrix has the form $\alpha I$, with $\alpha\in F^\times$, and it lies in $\SL(2,F)$ precisely when
\[
    \det(\alpha I)=\alpha^2=1.
\]
Thus $\alpha$ is a root of $X^2-1$. If $2\nmid q$, then $F$ has characteristic not equal to $2$, so
\[
    X^2-1=(X-1)(X+1)
\]
has the two distinct roots $1$ and $-1$. Hence $\lvert Z\cap \SL(2,q)\rvert=2$ in this case. If $2\mid q$, then $F$ has characteristic $2$, so $-1=1$ and
\[
    X^2-1=(X-1)^2
\]
has only the single root $1$. Hence $\lvert Z\cap \SL(2,q)\rvert=1$ in this case. Dividing the order of $\SL(2,q)$ by these two possible central orders gives the stated formula.
\end{proof}

\subsubsection*{Simplicity of $\mathrm{PSL}(n,F)$}
\begin{itemize}
    \item Our main objective in this section is to show that $\mathrm{PSL}(n,F)$ is simple whenever $n \geq 2$, except when $n = 2$ and $|F| = 2$ or $3$.
    \item This result dates back to L.~E.~Dickson, another of the early architects of the mathematical tradition of the University of Chicago, who established it in his 1896 Ph.D.\ thesis for $F$ a finite field; in 1893 E.~H.~Moore had established the result for $F$ a finite field and $n = 2$, while in 1870 Camille Jordan had established the result for $F$ a finite field of prime order and $n$ arbitrary.
    \item We first need two lemmas.
\end{itemize}

\begin{lemma}
\label{lem:6.6}
If $n \geq 2$, then every transvection $X_{ij}(\alpha)$ is a commutator of elements of $\mathrm{SL}(n,F)$, except when $n = 2$ and $|F| = 2$ or $3$.
\end{lemma}
\begin{proof}
First suppose that $n>2$. Given distinct $i$ and $j$, choose $k$ distinct from both of them. By Lemma~\ref{lem:4.6}(iv),
\[
    [X_{ik}(\alpha),X_{kj}(1)]=X_{ij}(\alpha).
\]
Both $X_{ik}(\alpha)$ and $X_{kj}(1)$ lie in $\SL(n,F)$ by Lemma~\ref{lem:4.6}(i), so $X_{ij}(\alpha)$ is a commutator of elements of $\SL(n,F)$.\\
It remains to consider the case $n=2$. Assume that $F$ is not one of the fields with $2$ or $3$ elements. Then there is some $\beta\in F^\times$ such that $\beta^2\neq 1$; otherwise every element of $F^\times$ would be a root of $X^2-1$, forcing $\lvert F^\times\rvert\leq 2$. Put
\[
    d=
    \begin{pmatrix}
        \beta & 0\\
        0     & \beta^{-1}
    \end{pmatrix}.
\]
Then $d\in\SL(2,F)$. For any $\gamma\in F$, a direct calculation gives
\[
    dX_{12}(\gamma)d^{-1}=X_{12}(\beta^2\gamma),
\]
and hence
\[
    [d,X_{12}(\gamma)]
    =X_{12}(\beta^2\gamma)X_{12}(-\gamma)
    =X_{12}((\beta^2-1)\gamma).
\]
Since $\beta^2-1\neq 0$, choosing $\gamma=\alpha(\beta^2-1)^{-1}$ gives
\[
    X_{12}(\alpha)=[d,X_{12}(\gamma)].
\]
Similarly,
\[
    dX_{21}(\gamma)d^{-1}=X_{21}(\beta^{-2}\gamma),
\]
so
\[
    [d,X_{21}(\gamma)]=X_{21}((\beta^{-2}-1)\gamma).
\]
Since $\beta^{-2}-1\neq0$, choosing $\gamma=\alpha(\beta^{-2}-1)^{-1}$ gives $X_{21}(\alpha)$ as a commutator of elements of $\SL(2,F)$. This proves the assertion in all non-exceptional cases.
\end{proof}

\subsubsection*{Consequence}
\begin{itemize}
    \item Observe that Lemma~\ref{lem:6.6} and Proposition~\ref{prop:6.2} together imply that, except when $n = 2$ and $|F| = 2$ or $3$, $\mathrm{SL}(n,F)$ is its own derived group and is consequently also the derived group of $\mathrm{GL}(n,F)$.
\end{itemize}

\begin{lemma}
\label{lem:6.7}
If $n \geq 2$, then the action of $\mathrm{SL}(n,F)$ on the one-dimensional subspaces of $V_n(F)$ is doubly transitive.
\end{lemma}
\begin{proof}
Let $U_1,U_2,W_1,W_2$ be one-dimensional subspaces of $V_n(F)$, with $U_1\neq U_2$ and $W_1\neq W_2$. Choose non-zero vectors
\[
    \mathbf{c}_1\in U_1,\qquad \mathbf{c}_2\in U_2,\qquad
    \mathbf{d}_1\in W_1,\qquad \mathbf{d}_2\in W_2.
\]
Since $U_1$ and $U_2$ are distinct one-dimensional subspaces, the vectors $\mathbf{c}_1$ and $\mathbf{c}_2$ are linearly independent. Similarly, $\mathbf{d}_1$ and $\mathbf{d}_2$ are linearly independent. Extend these pairs to bases
\[
    \mathbf{c}_1,\ldots,\mathbf{c}_n
    \qquad\text{and}\qquad
    \mathbf{d}_1,\ldots,\mathbf{d}_n
\]
of $V_n(F)$. Let $C$ and $D$ be the matrices whose $i$th columns are $\mathbf{c}_i$ and $\mathbf{d}_i$, respectively. Then $C,D\in \GL(n,F)$.\\
Put
\[
    \varepsilon=\frac{\det D}{\det C},
\]
and let $E$ be the diagonal matrix with first diagonal entry $\varepsilon$ and all other diagonal entries equal to $1$. Then
\[
    \det(DE^{-1}C^{-1})
    =(\det D)\varepsilon^{-1}(\det C)^{-1}
    =1,
\]
so
\[
    s=DE^{-1}C^{-1}\in \SL(n,F).
\]
Now $C^{-1}\mathbf{c}_1=\mathbf{v}_1$ and $C^{-1}\mathbf{c}_2=\mathbf{v}_2$, while $E^{-1}\mathbf{v}_1=\varepsilon^{-1}\mathbf{v}_1$ and $E^{-1}\mathbf{v}_2=\mathbf{v}_2$. Hence
\[
    s\mathbf{c}_1=\varepsilon^{-1}\mathbf{d}_1,
    \qquad
    s\mathbf{c}_2=\mathbf{d}_2.
\]
Thus $sU_1=W_1$ and $sU_2=W_2$. Therefore $\SL(n,F)$ sends any ordered pair of distinct one-dimensional subspaces to any other such ordered pair, so the action is doubly transitive.
\end{proof}

\begin{theorem}
\label{thm:6.8}
If $n \geq 2$, then $\mathrm{PSL}(n,F)$ is simple, except when $n = 2$ and $|F| = 2$ or $3$.
\end{theorem}
\begin{proof}
Assume that we are not in the exceptional cases. Put
\[
    S=\SL(n,F),
    \qquad
    L=F\mathbf{v}_1,
    \qquad
    Z_S=Z\cap S.
\]
It is enough to prove that if $N\trianglelefteq S$, then either $N\leq Z_S$ or $N=S$. Indeed, the normal subgroups of
\[
    \mathrm{PSL}(n,F)=S/Z_S
\]
correspond by \nameref{thm:correspondence} to the normal subgroups $N\trianglelefteq S$ with $Z_S\leq N\leq S$, via
\[
    N\longmapsto N/Z_S.
\]
If every normal subgroup of $S$ is either contained in $Z_S$ or equal to $S$, then any normal subgroup $N/Z_S$ of $S/Z_S$ has only two possibilities. In the first case, $Z_S\leq N\leq Z_S$, so $N=Z_S$, and hence
\[
    N/Z_S=Z_S/Z_S,
\]
which is the trivial subgroup of $S/Z_S$. In the second case, $N=S$, and hence
\[
    N/Z_S=S/Z_S,
\]
which is the whole quotient group. Thus proving the stated assertion about normal subgroups of $S$ proves that $\mathrm{PSL}(n,F)$ is simple.\\
Let $P$ be the stabilizer of the line $L$ under the action of $S$ on the one-dimensional subspaces of $V_n(F)$. By Lemma~\ref{lem:6.7} and Proposition~\ref{prop:3.8}, the subgroup $P$ is maximal in $S$. Now let
\[
    K=\left\{
    I+\sum_{j=2}^{n}\alpha_jE_{1j}
    \,\middle|\,
    \alpha_j\in F
    \right\}.
\]
Here $E_{1j}$ denotes the matrix with $1$ in the $(1,j)$-entry and $0$ everywhere else.
Equivalently, $K$ is the subgroup of upper unitriangular matrices whose only possible non-zero entries away from the main diagonal lie in the first row. Since $E_{1j}E_{1k}=0$ for all $j,k\geq2$, multiplication in $K$ just adds the parameters $\alpha_j$, and hence $K$ is abelian. Also, $K$ is exactly the subgroup of $P$ acting trivially on $L$ and on the quotient space $V_n(F)/L$; this description is invariant under conjugation by elements of $P$, because $P$ stabilizes $L$. Hence $K\trianglelefteq P$.\\
Let $N\trianglelefteq S$. First suppose that $N\leq P$. Then $N$ stabilizes $L$. If $s\in S$, then $sNs^{-1}=N$, and hence $N$ stabilizes $sL$. Since the action of $S$ on one-dimensional subspaces is transitive by Lemma~\ref{lem:6.7}, it follows that $N$ stabilizes every one-dimensional subspace of $V_n(F)$. In particular, every element of $N$ stabilizes each standard line $F\mathbf{v}_i$, so every element of $N$ is diagonal. If such an element has diagonal entries $\lambda_1,\ldots,\lambda_n$, then stabilizing each line $F(\mathbf{v}_i+\mathbf{v}_j)$ forces $\lambda_i=\lambda_j$ for all $i\neq j$. Thus every element of $N$ is scalar, and hence $N\leq Z_S$.\\
Now suppose that $N\nleq P$. Since $N\trianglelefteq S$, the product $PN$ is a subgroup of $S$. It properly contains $P$, so the maximality of $P$ gives
\[
    PN=S.
\]
Let $\eta\colon S\to S/N$ be the natural map. Since $K\trianglelefteq P$, we have $\eta(K)\trianglelefteq\eta(P)$. But
\[
    \eta(P)=PN/N=S/N,
\]
so $\eta(K)\trianglelefteq S/N$. By \nameref{thm:correspondence}, its inverse image
\[
    \eta^{-1}(\eta(K))=KN
\]
is normal in $S$. Since $K$ is generated by the root subgroups $X_{12},\ldots,X_{1n}$, and since $KN$ is normal in $S$, the subgroup $KN$ contains all $S$-conjugates of these root subgroups. By Proposition~\ref{prop:6.3}, all root subgroups are conjugate in $S$, so $KN$ contains every $X_{ij}$. Proposition~\ref{prop:6.2} then gives
\[
    KN=S.
\]
By \nameref{thm:first-iso},
\[
    S/N=KN/N\cong K/(K\cap N).
\]
Since $K$ is abelian, the quotient $S/N$ is abelian. Therefore Proposition~\ref{prop:2.6} gives $S'\leq N$. By Lemma~\ref{lem:6.6}, every transvection $X_{ij}(\alpha)$ is a commutator of elements of $S$, and hence every transvection lies in $N$. Since the transvections generate $S$ by Proposition~\ref{prop:6.2}, we conclude that $N=S$.\\
Thus every normal subgroup of $S$ is either contained in $Z_S$ or equal to $S$. Therefore $S/Z_S=\mathrm{PSL}(n,F)$ is simple.
\end{proof}

\subsection*{Exercises}

\begin{exercise}
\label{ex:6.1}
Let $B$ be the standard Borel subgroup of $\mathrm{GL}(n,F)$. Determine all subgroups of $\mathrm{SL}(n,F)$ which contain $B \cap \mathrm{SL}(n,F)$.
\end{exercise}
\begin{solution}
Let $S=\SL(n,F)$ and put $B_S=B\cap S$. We claim that the required subgroups are precisely
\[
    P\cap S,
\]
where $P$ ranges over the staircase groups of $\GL(n,F)$. Indeed, if $P$ is a staircase group, then $B\leq P$, and hence $B_S\leq P\cap S\leq S$.\\
Conversely, suppose that
\[
    B_S\leq H\leq S.
\]
If $n=1$, then $S=1$, so the assertion is immediate. Assume now that $n\geq2$. We first note that the only non-zero subspaces of $V_n(F)$ stabilized by $B_S$ are the standard flag subspaces $V_1,\ldots,V_n$. The proof is the same as Lemma~\ref{lem:5.3}, with one determinant correction. If
\[
    y=\alpha_1\mathbf{v}_1+\dots+\alpha_k\mathbf{v}_k\in V_k-V_{k-1},
\]
then $\alpha_k\neq0$, and because $n\geq2$ we can choose an upper triangular matrix of determinant $1$ whose $k$th column is $y$ by adjusting one diagonal entry different from the $k$th one. Thus the argument of Lemma~\ref{lem:5.3} applies with $B_S$ in place of $B$.\\
List the standard flag subspaces stabilized by $H$ as
\[
    V_{a_1},\ldots,V_{a_r},
    \qquad
    1\leq a_1<\dots<a_r=n.
\]
Let $P$ be the staircase group stabilizing the subflag
\[
    0\subset V_{a_1}\subset\dots\subset V_{a_r}=V_n(F).
\]
Then $H\leq P\cap S$. It remains to prove the reverse containment.\\
We use the same block argument as in Theorem~\ref{thm:5.4}. Fix a block
\[
    F\mathbf{v}_s\oplus\dots\oplus F\mathbf{v}_t,
    \qquad
    s=a_{k-1}+1,\quad t=a_k.
\]
If $s=t$, there is nothing to prove for this block. Assume $s<t$. Since there is no $H$-invariant standard flag subspace strictly between $V_{s-1}$ and $V_t$, there exists $h\in H$ such that the coefficient of $\mathbf{v}_t$ in $h\mathbf{v}_s$ is non-zero. Write $h=b_1wb_2$ by \nameref{thm:bruhat}, with $b_1,b_2\in B$ and $w$ a permutation matrix. Since $h\in S$, this Bruhat decomposition may be rewritten as
\[
    h=b_1'\dot{w}b_2',
\]
where $b_1',b_2'\in B_S$ and $\dot{w}\in S$ is a diagonal multiple of $w$. Indeed, choose a diagonal matrix $d$ with $\det(dw)=1$ and put $\dot{w}=dw$. If $\delta=\det(b_2)$, choose $t_0\in T$ with $\det(t_0)=\delta^{-1}$. Then
\[
    h=(b_1d^{-1}\dot{w}t_0^{-1}\dot{w}^{-1})\dot{w}(t_0b_2),
\]
and both outside factors lie in $B_S$. Since $h,b_1',b_2'\in H$, we get $\dot{w}\in H$.\\
The diagonal factor in $\dot{w}$ does not change the permutation pattern, so the same argument used in Theorem~\ref{thm:5.4} gives $w(s)=t$. Put $q=w^{-1}(s)$. Then $s<q\leq t$, so the upper root subgroup $X_{sq}$ lies in $B_S\leq H$. Conjugating by $\dot{w}$ sends this root subgroup onto $X_{ts}$, up to a non-zero scalar change in the parameter. Hence $X_{ts}\leq H$. Since all upper root subgroups inside the block already lie in $B_S$, the commutator identities of Lemma~\ref{lem:4.6}(iv) give every root subgroup $X_{ij}$ with $s\leq i\neq j\leq t$.\\
Therefore, in each block, $H$ contains all root subgroups supported inside that block. By Proposition~\ref{prop:6.2}, it contains the corresponding embedded special linear group on each block. It also contains every determinant-one diagonal matrix, since $T\cap S\leq B_S$. Hence $H$ contains the full group $L_P\cap S$, where $L_P$ is the standard Levi complement of $P$: any element of $L_P\cap S$ is block diagonal of total determinant $1$, and after multiplying by a suitable element of $T\cap S$ its blocks all have determinant $1$. Finally, the unipotent radical $U_P$ of $P$ lies in $B_S$, so Proposition~\ref{prop:5.2} gives
\[
    P\cap S=U_P(L_P\cap S)\leq H.
\]
Together with $H\leq P\cap S$, this proves $H=P\cap S$.
\end{solution}

\begin{exercise}
\label{ex:6.2}
Show that $\mathrm{PSL}(2,2) \cong S_3$ and $\mathrm{PSL}(2,3) \cong A_4$, and that these groups are not simple.
\end{exercise}
\begin{solution}
By Proposition~\ref{prop:6.5},
\[
    \lvert \mathrm{PSL}(2,2)\rvert=2^3-2=6
\]
and
\[
    \lvert \mathrm{PSL}(2,3)\rvert=\frac{3^3-3}{2}=12.
\]
First let $F$ be the field with two elements. The group $\SL(2,F)$ acts on the three non-zero vectors of $V_2(F)$. If a matrix fixes all three non-zero vectors, then in particular it fixes the basis vectors $\mathbf{v}_1$ and $\mathbf{v}_2$, so it is the identity. Hence this action gives an embedding
\[
    \mathrm{PSL}(2,2)=\SL(2,F)\hookrightarrow S_3.
\]
Since both groups have order $6$, we get $\mathrm{PSL}(2,2)\cong S_3$.\\
Now let $F$ be the field with three elements. The group $\SL(2,F)$ acts on the four one-dimensional subspaces of $V_2(F)$, namely
\[
    F\binom{1}{0},\quad
    F\binom{0}{1},\quad
    F\binom{1}{1},\quad
    F\binom{1}{-1}.
\]
The kernel of this action consists of the scalar matrices in $\SL(2,F)$: indeed, if a matrix fixes the first two lines, it is diagonal, and if it also fixes the line $F(\mathbf{v}_1+\mathbf{v}_2)$, then its two diagonal entries are equal. Thus the kernel is $\{I,-I\}$, and the action induces an embedding
\[
    \mathrm{PSL}(2,3)\hookrightarrow S_4.
\]
We show that the image lies in $A_4$. By Proposition~\ref{prop:6.2}, $\SL(2,F)$ is generated by $X_{12}$ and $X_{21}$. Write the four lines as
\[
    L_0=F\binom{1}{0},\quad
    L_1=F\binom{1}{1},\quad
    L_{-1}=F\binom{1}{-1},\quad
    M=F\binom{0}{1}.
\]
Then $X_{12}(1)$ acts as the cycle
\[
    (L_1\;L_{-1}\;M)
\]
and $X_{21}(1)$ acts as the cycle
\[
    (L_0\;L_1\;L_{-1}).
\]
These are even permutations, and their powers give the images of the whole root subgroups $X_{12}$ and $X_{21}$. Therefore the image of $\mathrm{PSL}(2,3)$ in $S_4$ is contained in $A_4$. Since this image has order $12$ and $\lvert A_4\rvert=12$, it follows that $\mathrm{PSL}(2,3)\cong A_4$.\\
Finally, these groups are not simple: $A_3\trianglelefteq S_3$ is a non-trivial proper normal subgroup, and the Klein four-subgroup is a non-trivial proper normal subgroup of $A_4$.
\end{solution}

\begin{exercise}
\label{ex:6.3}
Show that $\mathrm{PSL}(2,4)$ and $\mathrm{PSL}(2,5)$ are both isomorphic with $A_5$.
\end{exercise}
\begin{solution}
By Proposition~\ref{prop:6.5},
\[
    \lvert \mathrm{PSL}(2,4)\rvert=4^3-4=60
\]
and
\[
    \lvert \mathrm{PSL}(2,5)\rvert=\frac{5^3-5}{2}=60.
\]
First let $F$ be the field with four elements. The group $\SL(2,F)$ acts on the five one-dimensional subspaces of $V_2(F)$. The kernel of this action consists of the scalar matrices in $\SL(2,F)$, and since $F$ has characteristic $2$, the only scalar matrix in $\SL(2,F)$ is $I$. Thus this action gives an embedding
\[
    \mathrm{PSL}(2,4)=\SL(2,F)\hookrightarrow S_5.
\]
We show that the image lies in $A_5$. Write the five lines as
\[
    L_a=F\binom{1}{a}\quad(a\in F),
    \qquad
    M=F\binom{0}{1}.
\]
For $\alpha\neq0$, the transvection $X_{21}(\alpha)$ fixes $M$ and sends
\[
    L_a\mapsto L_{a+\alpha}.
\]
Since the additive group of $F$ has characteristic $2$, the translation $a\mapsto a+\alpha$ on the four elements of $F$ is a product of two disjoint transpositions. Hence $X_{21}(\alpha)$ acts as an even permutation on the five lines. Similarly, $X_{12}(\alpha)$ fixes $L_0$ and, for $\alpha\neq0$, fixes no other line: if
\[
    X_{12}(\alpha)\binom{1}{a}
    =
    \binom{1+\alpha a}{a}
\]
spans the same line as $\binom{1}{a}$ and $a\neq0$, then the second coordinate forces the scalar to be $1$, and hence $\alpha a=0$, a contradiction. Also $M$ is sent to $L_{\alpha^{-1}}$. Thus $X_{12}(\alpha)$ is an involution fixing exactly one of the five lines, so it is again a product of two disjoint transpositions. By Proposition~\ref{prop:6.2}, the image of $\SL(2,F)$ in $S_5$ is generated by these even permutations, and hence lies in $A_5$. Since the image has order $60=\lvert A_5\rvert$, we conclude that
\[
    \mathrm{PSL}(2,4)\cong A_5.
\]
Now let $F$ be the field with five elements, and put $G=\mathrm{PSL}(2,5)$. We shall find a subgroup of $G$ of index $5$. In $\SL(2,F)$, let
\[
    C=
    \begin{pmatrix}
        0  & 1\\
        -1 & 0
    \end{pmatrix},
    \qquad
    D=
    \begin{pmatrix}
        2 & 0\\
        0 & -2
    \end{pmatrix},
    \qquad
    R=
    \begin{pmatrix}
        1 & 1\\
        2 & 3
    \end{pmatrix}.
\]
Let $c,d,r$ be their images in $G$. Since
\[
    C^2=D^2=(CD)^2=-I
\]
and $CD=-DC$, the subgroup
\[
    V=\{1,c,d,cd\}
\]
is a Klein four-subgroup of $G$. Also $R^3=I$, so $r$ has order $3$. A direct calculation, modulo the scalar subgroup $\{I,-I\}$, gives
\[
    rcr^{-1}=cd,\qquad
    r(cd)r^{-1}=d,\qquad
    rdr^{-1}=c.
\]
Thus $r$ normalizes $V$ and cyclically permutes its three non-identity elements. Therefore
\[
    H=V\langle r\rangle
\]
is a subgroup of $G$ of order $12$. Since $\lvert G\rvert=60$, the index of $H$ in $G$ is $5$.\\
Let $G$ act on the five left cosets of $H$. This gives a homomorphism
\[
    G\to S_5.
\]
The kernel is normal in $G$. By Theorem~\ref{thm:6.8}, the group $G=\mathrm{PSL}(2,5)$ is simple, since this is not one of the exceptional cases. The action is not trivial because $H$ is a proper subgroup, so the kernel cannot be all of $G$. Hence the kernel is trivial, and $G$ embeds in $S_5$. Finally, composing this embedding with the sign homomorphism $S_5\to\{\pm1\}$ gives a homomorphism from the simple group $G$ to a group of order $2$. This homomorphism must be trivial, since a non-trivial one would have a non-trivial proper kernel. Therefore the image of $G$ lies in $A_5$. Since both $G$ and $A_5$ have order $60$, the image is all of $A_5$, and hence
\[
    \mathrm{PSL}(2,5)\cong A_5.
\]
\end{solution}

\begin{exercise}
\label{ex:6.4}
Show that $\mathrm{PSL}(2,7) \cong \mathrm{GL}(3,2)$.
\end{exercise}
\begin{solution}
Let $F$ be the field with seven elements, and put $G=\mathrm{PSL}(2,7)$. By Proposition~\ref{prop:6.5},
\[
    |G|=\frac{7^3-7}{2}=168.
\]
Also
\[
    |\mathrm{GL}(3,2)|=(2^3-1)(2^3-2)(2^3-2^2)=7\cdot 6\cdot 4=168.
\]
We shall construct the same faithful permutation group of degree $7$ from $G$ and from $\mathrm{GL}(3,2)$.\\
First we construct a subgroup of $G$ of index $7$. In $\SL(2,F)$, let
\[
    C=
    \begin{pmatrix}
        0  & 1\\
        -1 & 0
    \end{pmatrix},
    \quad
    D=
    \begin{pmatrix}
        2 & 3\\
        3 & -2
    \end{pmatrix},
    \quad
    R=
    \begin{pmatrix}
        0 & 2\\
        3 & 1
    \end{pmatrix},
    \quad
    E=
    \begin{pmatrix}
        1 & -1\\
        2 & -1
    \end{pmatrix}.
\]
Let $c,d,r,e$ be their images in $G$. A direct calculation gives
\[
    C^2=D^2=(CD)^2=-I,\qquad CD=-DC.
\]
Thus
\[
    V=\{1,c,d,cd\}
\]
is a Klein four-subgroup of $G$. Also $R^3=-I$ and $E^2=-I$, so $r$ has order $3$ and $e$ has order $2$ in $G$. Conjugation gives
\[
    rcr^{-1}=d,\qquad rdr^{-1}=cd,\qquad r(cd)r^{-1}=c,
\]
and
\[
    ece^{-1}=cd,\qquad ede^{-1}=d,\qquad e(cd)e^{-1}=c.
\]
Moreover $ere^{-1}=r^{-1}$, so $\langle r,e\rangle\cong S_3$. Hence $r$ and $e$ act on the three non-identity elements of $V$ as a $3$-cycle and a transposition. This action is faithful, while every element of $V$ centralizes $V$, so $V\cap \langle r,e\rangle=1$. Therefore the subgroup
\[
    H=V\langle r,e\rangle
\]
has order $4\cdot 6=24$. Since $|G|=168$, the subgroup $H$ has index $7$ in $G$.\\
Let $G$ act on the seven left cosets of $H$. The kernel is normal in $G$. By Theorem~\ref{thm:6.8}, the group $G=\mathrm{PSL}(2,7)$ is simple. The action is not trivial because $H$ is proper, so the kernel is not all of $G$. Hence the kernel is trivial, and this coset action embeds $G$ in $S_7$.\\
Now take
\[
    S=
    \begin{pmatrix}
        0  & -1\\
        1  & 0
    \end{pmatrix},
    \qquad
    T=
    \begin{pmatrix}
        1 & 1\\
        0 & 1
    \end{pmatrix}.
\]
The images of $S$ and $T$ generate $G$: indeed, $T=X_{12}(1)$, and conjugating powers of $T$ by $S$ gives the lower root subgroup $X_{21}$, so Proposition~\ref{prop:6.2} applies. We now make the coset calculation explicit. Choose the following representatives $g_i$ for the seven left cosets of $H$ in $G$, where each displayed matrix denotes its image in $G$:
\[
\begin{array}{c|c}
    i & g_i\\
    \hline
    1 & \bigl(\begin{smallmatrix}2&-3\\2&1\end{smallmatrix}\bigr)\\
    2 & \bigl(\begin{smallmatrix}1&0\\2&1\end{smallmatrix}\bigr)\\
    3 & \bigl(\begin{smallmatrix}2&3\\-2&1\end{smallmatrix}\bigr)\\
    4 & \bigl(\begin{smallmatrix}3&0\\2&-2\end{smallmatrix}\bigr)\\
    5 & \bigl(\begin{smallmatrix}3&-2\\0&-2\end{smallmatrix}\bigr)\\
    6 & \bigl(\begin{smallmatrix}3&2\\0&-2\end{smallmatrix}\bigr)\\
    7 & \bigl(\begin{smallmatrix}1&-1\\2&-1\end{smallmatrix}\bigr).
\end{array}
\]
For $X\in\{S,T\}$, the equality $Xg_iH=g_jH$ is equivalent to
\[
    g_j^{-1}Xg_i\in H.
\]
Checking this membership for the displayed representatives gives
\[
\begin{array}{c|ccccccc}
    i & 1&2&3&4&5&6&7\\
    \hline
    Sg_iH & g_1H&g_5H&g_4H&g_3H&g_2H&g_6H&g_7H\\
    Tg_iH & g_3H&g_5H&g_6H&g_7H&g_4H&g_2H&g_1H.
\end{array}
\]
For example,
\[
    g_5^{-1}Sg_2=
    \begin{pmatrix}
        1  & -2\\
        -3 & 0
    \end{pmatrix}
    =R^{-1}\in H,
\]
so $Sg_2H=g_5H$. Reading the two rows of the table as permutations of the labels $1,\ldots,7$, we get
\[
    S\mapsto (2\;5)(3\;4),
    \qquad
    T\mapsto (1\;3\;6\;2\;5\;4\;7).
\]
Now consider the natural action of $\mathrm{GL}(3,2)$ on the seven non-zero vectors of $\mathbb F_2^3$. Label these vectors by
\[
\begin{array}{c|ccccccc}
    \text{label} & 1&2&3&4&5&6&7\\
    \hline
    \text{vector} & e_1&e_2&e_1+e_2&e_3&e_1+e_3&e_2+e_3&e_1+e_2+e_3.
\end{array}
\]
With respect to the basis $e_1,e_2,e_3$, the matrices
\[
    A=
    \begin{pmatrix}
        1&1&1\\
        0&0&1\\
        0&1&0
    \end{pmatrix},
    \qquad
    B=
    \begin{pmatrix}
        1&1&1\\
        1&0&1\\
        0&1&1
    \end{pmatrix}
\]
are invertible over $\mathbb F_2$. A direct calculation on the seven labelled vectors gives
\[
    A\mapsto (2\;5)(3\;4),
    \qquad
    B\mapsto (1\;3\;6\;2\;5\;4\;7).
\]
Thus the two generators of the image of $G$ in $S_7$ lie in the image of the natural embedding
\[
    \mathrm{GL}(3,2)\hookrightarrow S_7.
\]
Since $S$ and $T$ generate $G$, the whole image of $G$ lies in the image of $\mathrm{GL}(3,2)$. Both images have order $168$: the coset action of $G$ is faithful, and the natural action of $\mathrm{GL}(3,2)$ on the non-zero vectors is faithful. Therefore the two images in $S_7$ are equal, and hence
\[
    \mathrm{PSL}(2,7)\cong \mathrm{GL}(3,2).
\]
\end{solution}

\begin{exercise}
\label{ex:6.5}
Show that $\mathrm{PSL}(2,9) \cong A_6$.
\end{exercise}
\begin{solution}
\end{solution}

\begin{exercise}
\label{ex:6.6}
By Exercise~\ref{ex:6.3} and Theorem~\ref{thm:6.8}, the group $A_5$ is simple. Attempt to prove this fact directly by mimicking the proof of Theorem~\ref{thm:6.8}, with $A_5$ in place of $S$, $A_4$ in place of $P$, and the Klein four-group in place of $K$.
\end{exercise}
\begin{solution}
\end{solution}

\subsection*{Further Exercises}

Let $F$ be a field, and let $V = F^2$. We define an equivalence relation on $V - \{0\}$ by $\mathbf{v} \sim \mathbf{w}$ iff $\mathbf{v} = \alpha \mathbf{w}$ for some $\alpha \in F^{\times}$. The set of equivalence classes under this relation is called the \emph{one-dimensional projective space} (or \emph{projective line}) over $F$ and is denoted by $\mathbb{P}^1(F)$. We write elements of $V$ as column vectors, and the element of $\mathbb{P}^1(F)$ that is the equivalence class of $\bigl(\begin{smallmatrix}x\\ y\end{smallmatrix}\bigr)$ shall be written as $\bigl[\begin{smallmatrix}x\\ y\end{smallmatrix}\bigr]$.

\begin{exercise}
\label{ex:6.7}
Show that there is a well-defined action of $\mathrm{PSL}(2,F)$ on $\mathbb{P}^1(F)$ given by
\[
\overline{\begin{pmatrix}a & b\\ c & d\end{pmatrix}} \begin{bmatrix}x\\ y\end{bmatrix} = \begin{bmatrix}ax+by\\ cx+dy\end{bmatrix}
\]
and that this action is doubly transitive. (Here $\overline{\bigl(\begin{smallmatrix}a & b\\ c & d\end{smallmatrix}\bigr)}$ denotes the image in $\mathrm{PSL}(2,F)$ of $\bigl(\begin{smallmatrix}a & b\\ c & d\end{smallmatrix}\bigr) \in \mathrm{SL}(2,F)$.)
\end{exercise}
\begin{solution}
\end{solution}

\begin{exercise}
\label{ex:6.8}
(cont.) Determine the proper definition of $n$-dimensional projective space $\mathbb{P}^n(F)$ for arbitrary $n$, and show that $\mathrm{PSL}(n+1,F)$ acts doubly transitively on $\mathbb{P}^n(F)$.
\end{exercise}
\begin{solution}
\end{solution}

\begin{exercise}
\label{ex:6.9}
Let $F$ be the field of $7$ elements. For $0 \leq i \leq 6$, let $i$ represent $\bigl[\begin{smallmatrix}1\\ i\end{smallmatrix}\bigr] \in \mathbb{P}^1(F)$; denote the remaining element $\bigl[\begin{smallmatrix}0\\ 1\end{smallmatrix}\bigr]$ of $\mathbb{P}^1(F)$ by $\infty$. View $S_8$ as being the group of permutations of $\{\infty, 0, 1, 2, 3, 4, 5, 6\}$, and let $\varphi\colon \mathrm{PSL}(2,7) \to S_8$ be the monomorphism arising (as in Proposition~\ref{prop:3.1}) from the action of $\mathrm{PSL}(2,7)$ on $\mathbb{P}^1(F)$. Show that the image of $\varphi$ is the subgroup of $A_8$ generated by $(0\;1\;2\;3\;4\;5\;6)$, $(1\;4\;2)(3\;5\;6)$, and $(\infty\;0)(1\;3)(2\;5)(4\;6)$. (\textsc{Hint:} Show first that $\mathrm{PSL}(2,7)$ is generated by the images under the natural map of the elements $\bigl(\begin{smallmatrix}1 & 0\\ 1 & 1\end{smallmatrix}\bigr)$, $\bigl(\begin{smallmatrix}4 & 0\\ 0 & 2\end{smallmatrix}\bigr)$, and $\bigl(\begin{smallmatrix}0 & 3\\ 2 & 0\end{smallmatrix}\bigr)$ of $\mathrm{SL}(2,7)$.) For the relevance of this exercise, see Exercises~\ref{ex:7.13}--\ref{ex:7.17}.
\end{exercise}
\begin{solution}
\end{solution}

Recall the definition of a $BN$-pair from the further exercises to Section~4.

\begin{exercise}
\label{ex:6.10}
Let $G = \mathrm{GL}(n,F)$. Let $B$ be the standard Borel subgroup of $G$, let $N$ be the subgroup of $G$ of monomial matrices, and let $T = B \cap N$ be the subgroup of diagonal matrices. Let $B_0 = B \cap \mathrm{SL}(n,F)$, let $N_0 = N \cap \mathrm{SL}(n,F)$, and let $T_0 = T \cap \mathrm{SL}(n,F) = B_0 \cap N_0$. Show that $B_0$ and $N_0$ form a $BN$-pair of $\mathrm{SL}(n,F)$, with the associated Weyl group $W_0 = N_0/T_0$ being isomorphic with $S_n$.
\end{exercise}
\begin{solution}
\end{solution}

\begin{exercise}
\label{ex:6.11}
Let $P = \mathrm{PSL}(n,F)$, and let $\overline{B_0}$, $\overline{N_0}$, and $\overline{T_0}$ be the images in $P$ of the subgroups $B_0$, $N_0$, and $T_0$ of $\mathrm{SL}(n,F)$. Show that $\overline{B_0}$ and $\overline{N_0}$ form a $BN$-pair of $P$ (with again the Weyl group $\overline{W_0} = \overline{N_0}/\overline{T_0}$ being isomorphic with $S_n$).
\end{exercise}
\begin{solution}
\end{solution}

\end{document}
